%mac
\immini{
\subsubsection{Sự rơi của các vật trong không khí}
\newghichu{Các vật rơi trong không khí xảy ra nhanh chậm khác nhau là do lực cản của không khí tác dụng vào chúng khác nhau.}
\subsubsection{Sự rơi của các vật trong chân không (sự rơi tự do)}
\newghichu{Sự rơi tự do là sự rơi của một vật chỉ chịu tác dụng của trọng lực}
\subsubsection{Đặc điểm của sự rơi tự do}
\begin{itemize}
\item Sự rơi tự do có phương thẳng đứng (phương của dây rọi), có chiều chuyển động từ trên xuống.
\item Sự rơi tự do là chuyển động thẳng nhanh dần đều với gia tốc \bth{a=g=const}.
\item Tại cùng một nơi trên Trái Đất và ở gần mặt đất, các vật đều rơi tự do với cùng gia tốc g.
\end{itemize}

}
{\begin{tikzpicture}[scale=1,thick,>=stealth']
\node at (0,0.1){\includegraphics[width=0.45cm]{images/LongChim}};
\node at (0.4,0){\includegraphics[width=0.45cm]{images/Bi}};
\draw[blue,fill=white,opacity=0.3](-0.35,-1)--++(0:1)--++(90:6)coordinate(a)--++(180:1)--cycle;
\draw[blue,fill=white,opacity=0.3](a)--++(90:0.3)--++(180:1)--++(270:0.3)--++(0:0.4)coordinate(b)--cycle;
\draw[blue,fill=gray,opacity=0.7](b)--++(90:0.5)--++(0:0.2)--++(-90:0.5)--cycle;
\draw(0.2,2)--++(30:1)node[above right]{Chân}node[below right]{không};
\begin{scope}[shift={(-3,0)}]
\node at (0.1,3){\includegraphics[width=0.45cm]{images/LongChim}};
\node at (0.1,0){\includegraphics[width=0.45cm]{images/Bi}};
\draw[blue,fill=white,opacity=0.3](-0.35,-1)--++(0:1)--++(90:6)coordinate(a)--++(180:1)--cycle;
\draw[blue,fill=white,opacity=0.3](a)--++(90:0.3)--++(180:1)--++(270:0.3)--++(0:0.4)coordinate(b)--cycle;
\draw[blue,fill=gray,opacity=0.7](b)--++(90:0.5)--++(0:0.2)--++(-90:0.5)--cycle;
\draw(0.2,2)--++(30:1)node[above right]{Không}node[below right]{khí};
\end{scope}
\end{tikzpicture}}
\begin{note}
Nếu không đòi hỏi độ chính xác cao ta thường lấy \bth{g=9,8\ m/s^2} hoặc \bth{g=10\ m/s^2}.
\end{note}
\subsubsection{Các công thức của rơi tự do}
\immini{\begin{itemize}
\item Vận tốc: $v=gt$\\
\item Độ dịch chuyển - Quãng đường: $d=s=\dfrac{gt^2}{2}$\\
\item Công thức liên hệ: \ \ \ \ \ \ $v^2=2gs$
\end{itemize}}{\begin{tikzpicture}[scale=1.1,thick,>=stealth',draw=Mapcolor]
\tkzDefPoints{0/0/o, -0.1/-2/g, 0/-2/a, 0/-5/y, 0.2/0/s, 0.1/-2/v, -1/-5/m, 1/-5/n, -1/-5.15/q, 1/-5.15/p}
\tkzDrawPoints[size=2](o,a,g,v)
\fill[pattern=north east lines](m)--(n)--(p)--(q);
\tkzDrawSegments[dashed](o,a)
\tkzDrawSegments(m,n)
\tkzDrawSegments[->](a,y)
\draw[very thick,Mapcolor,->](v)--++(-90:2)node[right]{$\vec{v}$};
\draw[very thick,Mapcolor,->](g)--++(-90:1)node[left]{$\vec{g}$};
\tkzLabelPoint[left](o){$O$}
\tkzLabelPoint[above left](y){}
\begin{scope}[/pgf/decoration/raise=5pt]
\draw [decorate,decoration={brace,mirror,amplitude=5pt},xshift=0pt,yshift=-4pt]
(a) -- (o) node [black,midway,xshift=20pt] 
{$s$};
\end{scope}
\end{tikzpicture}}